\documentclass[11pt]{article}
\usepackage[margin=1in]{geometry}
\usepackage{amsmath,amssymb,amsthm,mathtools}
\IfFileExists{lmodern.sty}{\usepackage{lmodern}}{}
\usepackage[expansion=false]{microtype}
\usepackage{booktabs,array}
\usepackage{enumitem}
\usepackage{xcolor}
\usepackage[colorlinks=true,linkcolor=blue!50!black,citecolor=green!35!black,urlcolor=blue!50!black]{hyperref}
\usepackage[nameinlink]{cleveref}
\newcolumntype{R}[1]{>{\raggedright\arraybackslash}p{#1}}

\newtheorem{theorem}{Theorem}[section]
\newtheorem{proposition}[theorem]{Proposition}
\newtheorem{lemma}[theorem]{Lemma}
\newtheorem{corollary}[theorem]{Corollary}
\newtheorem{conjecture}[theorem]{Conjecture}
\theoremstyle{definition}
\newtheorem{definition}[theorem]{Definition}
\newtheorem{example}[theorem]{Example}
\newtheorem{question}[theorem]{Question}
\theoremstyle{remark}
\newtheorem{remark}[theorem]{Remark}

\newcommand{\D}{\mathsf{D}}
\newcommand{\Kl}{\mathrm{Kl}}
\newcommand{\Cpl}{\operatorname{Cpl}}
\newcommand{\TV}{\mathrm{TV}}
\newcommand{\Hyp}{\mathrm{Hyp}}
\newcommand{\Sh}{\mathrm{Sh}}
\newcommand{\C}{\mathsf{C}}
\newcommand{\supp}{\operatorname{supp}}
\newcommand{\id}{\mathrm{id}}

\title{Conditionals, Couplings and Toric Geometry:\\
Established Results and a Fukaya Programme\\[4pt]
\large Revised draft}
\author{Vinay K. Awasthi\\
Department of Applied Mathematics and Statistics\\
Johns Hopkins University\thanks{Verification scripts: \texttt{check\_glue.py},
\texttt{check\_stochastic\_basechange.py}, \texttt{check\_conditionals.py}; see \cref{sec:code}.}}
\date{October 4, 2026}

\begin{document}
\maketitle

\begin{abstract}
We separate the conditional/coupling formulation of the finite distribution monad into what is proved,
what is imported from symplectic topology, and what is conjectural. The proved layer has four parts.
(1) Markov kernels are the 1-cells of the Kleisli category; couplings are 2-cells whose vertical
composition is the gluing lemma; the natural whiskerings are not functorial, and the strict
2-categorical structure is a graded metric shadow. (2) The base-relative posterior tower is a fibre
over this base: its coherence, canonical section, strict mixed distributive law, admissibility data
and obstructions are all transported along Markov kernels. (3) Coupling polytopes are moment
polytopes: for nondegenerate marginals the polytope of couplings is Delzant, so it is the image of the
moment map of a compact toric symplectic manifold, and when every cell is a facet the Fisher metric on
the polytope is twice the Hessian of Guillemin's canonical symplectic potential. (4) The marginal
space is cut by walls $\sum_If=\sum_Jg$; the combinatorial type, and hence the toric manifold, is
locally constant off the walls and changes across them; the manifold is monotone exactly for uniform
marginals, which are nondegenerate exactly for coprime sizes, and the smallest example is the degree-6
del Pezzo surface. The imported layer uses known Floer theory to locate a distinguished Lagrangian
torus fibre in the monotone case over the \emph{product} coupling, not the maximal coupling. The
conjectural layer is restated with correct types: a family of toric manifolds over the chamber complex
with wall-crossing in place of monodromy, and the whiskering defect as the second component of an
$A_\infty$-functor rather than an $m_3$. We also list the corrections to the first draft.
\end{abstract}

\tableofcontents

%=====================================================================
\section{Introduction}\label{sec:intro}
%=====================================================================

Two structures sit on top of the finite distribution monad $\D$. The base-relative posterior tower
\cite{Awasthi-Tower} fixes an outcome base $B$ and solves a \emph{strict integrability} problem: it
carries a mixed distributive law satisfying all four Beck axioms. The conditional/coupling formulation
\cite{Awasthi-Attribution} lets the base vary along Markov kernels and compares kernels by couplings;
it solves a \emph{transport} problem. The relation is
\begin{align*}
\text{base-relative tower}&=\text{fibre over a fixed base},\\
\text{conditional/coupling structure}&=\text{the varying-base global structure}.
\end{align*}
The first draft of this note left one gap: the tower's change of base had been established only for
functions $f\colon B\to B'$, while the global structure uses kernels $k\colon B\to\D(B')$. That gap is
closed in \cite[\S9.7]{Awasthi-Tower}; we state the result in \cref{sec:fibre}.

The second purpose of this note is the geometry. The first draft conjectured a symplectic structure
whose restriction to a fibre ``is the Fisher metric'', and that the critical points of a projection are
the maximal couplings. The correct statements turn out to be theorems about coupling polytopes
(\cref{sec:toric}), and they point in a different direction: the distinguished point is the product
coupling.

\paragraph{Status of the results.}
\begin{center}\small
\begin{tabular}{@{}R{0.2\textwidth}R{0.48\textwidth}R{0.22\textwidth}@{}}
\toprule
Layer & Content & Where\\
\midrule
Proved & kernels, coupling 2-cells, graded shadow; stochastic base change of the tower; coupling
polytopes are Delzant; Fisher $=2\,\mathrm{Hess}$ of Guillemin's potential; walls and chambers;
monotone cases & \cref{sec:established,sec:fibre,sec:toric}\\
Imported & Lagrangian torus fibres; nondisplaceability and Floer cohomology of the special fibre in the
monotone case & \cref{sec:imported}\\
Conjectural & families over the chamber complex; whiskering defects as $A_\infty$-functor data; a
Fukaya-type functor & \cref{sec:conjectural}\\
\bottomrule
\end{tabular}
\end{center}

%=====================================================================
\section{The established layer: kernels and couplings}\label{sec:established}
%=====================================================================

$\D$ is the finitely supported distribution monad on $\mathbf{Set}$ and $\Kl(\D)$ its Kleisli category.
A 1-cell $k\colon B\rightsquigarrow B'$ is a Markov kernel $k\colon B\to\D(B')$, acting on states by
$k^\#(p)(b')=\sum_bp(b)k(b)(b')$. The results of this section are proved in
\cite[\S7]{Awasthi-Attribution}; we restate them for reference.

\begin{proposition}[1-cells {\cite[Props.~7.2, 7.3]{Awasthi-Attribution}}]\label{prop:1cells}
$k^\#$ is affine, $(h\circ k)^\#=h^\#k^\#$, $k^\#\circ\mu=\mu\circ\D(k^\#)$, and
$\TV(k^\#p,k^\#q)\le\TV(p,q)$. Attribution sets are transported along $k$, and precision can only be lost.
\end{proposition}

\begin{definition}[Coupling 2-cells]\label{def:coupling}
For 1-cells $f,g\colon X\rightsquigarrow Y$, a coupling 2-cell $\alpha\colon f\Rightarrow g$ is a kernel
$\alpha\colon X\to\D(Y\times Y)$ with marginals $f(x)$ and $g(x)$. Its cost is
$c(\alpha)=\max_x\Pr_{\alpha(x)}(y\ne y')$; the identity is the diagonal coupling; the vertical
composite glues along the common marginal,
$(\beta\cdot\alpha)(x)(y,y'')=\sum_{y'}\alpha(x)(y,y')\beta(x)(y',y'')/g(x)(y')$.
\end{definition}

\begin{theorem}[Witnesses and shadow {\cite[Thms.~7.6, 7.8]{Awasthi-Attribution}}]\label{thm:witness}
\leavevmode
\begin{enumerate}[label=(\roman*)]
\item For fixed $X,Y$, 1-cells and coupling 2-cells form a category under gluing, strictly
  associative and unital. Cost is subadditive and $\min_\alpha c(\alpha)=d(f,g)=\max_x\TV(f(x),g(x))$.
\item Whiskering by a deterministic precomposition is strictly functorial. Whiskering by a stochastic
  precomposition (mixing couplings) or by a postcomposition (copying on equal pairs, independent on
  unequal pairs) preserves identities, marginals and the cost bound, but is not functorial.
\item Graded 2-cells $f\Rightarrow_\varepsilon g$, $\varepsilon\ge d(f,g)$, with grades adding, form a strict
  2-category on kernel 1-cells, and $d$ makes $\Kl(\D)$ a Lawvere-enriched category.
\end{enumerate}
\end{theorem}

%=====================================================================
\section{The fibre: stochastic change of base for the tower}\label{sec:fibre}
%=====================================================================

The base-relative tower over $B$ has context algebra $A_B=\D B\times\D^2B\times\D^3B\times\Sh(B)$,
reader comonad $\C^B(Y)=Y\times A_B$ and strict law $\lambda^B$ \cite[\S9]{Awasthi-Tower}. With
$\Hyp=\Theta\times\D(-)$, a kernel $k\colon B\to\D(B')$ acts on hypotheses by
$\Hyp(k)(\theta,\rho)=(\theta,k^\#\rho)$, and we put
$A_k=k^\#\times\D(k^\#)\times\D^2(k^\#)\times\D(\Hyp k)$.

\begin{theorem}[The tower is a functor on the Kleisli category {\cite[Prop.~9.14, Cor.~9.15, Prop.~9.16]{Awasthi-Tower}}]\label{thm:fibre}
\leavevmode
\begin{enumerate}[label=(\roman*)]
\item $A_k$ is a morphism of $\D$-algebras, preserves coherent contexts, commutes with the canonical
  section $\sigma$, and $A_{h\circ k}=A_hA_k$; for deterministic $k$ it is the earlier base change.
\item $\id\times A_k$ is a comonad morphism $\C^B\Rightarrow\C^{B'}$ commuting with the laws $\lambda^B$,
  $\lambda^{B'}$.
\item Admissibility data, coherent lifts and the effective sets are transported; certificates pull
  back along $f'\mapsto k\rhd f'$; failure levels can only rise and the secondary obstruction can only
  fall.
\item Conditioning on an outcome observed after the channel equals conditioning on the pulled-back
  predicate $k\rhd1_{e'}$, then changing base.
\end{enumerate}
Graded 2-cells transport to every level: for kernels $f,g$, the level-one and level-two Kantorovich
distances between $A_f(c)$ and $A_g(c)$ are at most $d(f,g)$.
\end{theorem}

The proof uses a single fact beyond naturality: Kleisli extension is a morphism of free $\D$-algebras.
All other components of $A_k$ are pushforwards along functions. So the strict integrability of the
fibre and the transport of the global structure are compatible, and the fibre is functorial on all
of $\Kl(\D)$, not only on $\mathbf{Set}$.

%=====================================================================
\section{Coupling polytopes are moment polytopes}\label{sec:toric}
%=====================================================================

\subsection{The polytope and its lattice}

Let $f\in\D(Y)$ and $g\in\D(Y')$ have full support, with $\lvert Y\rvert=n$, $\lvert Y'\rvert=m$. (For
coupling 2-cells $Y=Y'$; rectangular polytopes arise when gluing contexts
\cite[Prop.~7.11]{Awasthi-Attribution}.) The \emph{coupling polytope} is
\[
\Cpl(f,g)=\Big\{\alpha\in\mathbb{R}_{\ge0}^{Y\times Y'}:\ \textstyle\sum_{y'}\alpha(y,y')=f(y),\
\sum_y\alpha(y,y')=g(y')\Big\}.
\]
Its tangent space is the space of \emph{circulations}, matrices with zero row and column sums. The
integer circulations form a lattice $\Lambda$ of rank $(n-1)(m-1)$; a basis is
$E_{ij}-E_{i,m}-E_{n,j}+E_{n,m}$ for $i<n$, $j<m$.

\begin{proposition}[Interior and Fisher metric]\label{prop:interior}
$\Cpl(f,g)$ has dimension $(n-1)(m-1)$, its relative interior is the set of strictly positive
couplings, and on it the Fisher metric $\mathcal{I}_\alpha(v,w)=\sum v(y,y')w(y,y')/\alpha(y,y')$ is
positive definite on circulations.
\end{proposition}
\begin{proof}
The product coupling $f\otimes g$ is strictly positive, so no inequality $\alpha\ge0$ is an implicit
equality, and the relative interior of the polytope in its affine hull is $\{\alpha>0\}$. There
$\mathcal{I}_\alpha(v,v)=\sum v^2/\alpha>0$ for $v\ne0$.
\end{proof}

\begin{remark}
The first draft argued nondegeneracy through a radical and then distinguished ``the Fisher metric''
from ``the Fisher--Rao metric''. These are the same metric, the Fisher information metric. The
square-root embedding is a way of computing its geodesic distance \cite[Prop.~8.1]{Awasthi-Attribution}.
\end{remark}

\subsection{Nondegeneracy implies Delzant}

\begin{definition}[Walls]\label{def:walls}
The pair $(f,g)$ is \emph{degenerate} if $\sum_{y\in I}f(y)=\sum_{y'\in J}g(y')$ for some
nonempty proper subsets $I\subsetneq Y$ and $J\subsetneq Y'$. These equations define hyperplanes,
the \emph{walls}, in the space of pairs of marginals; the \emph{chambers} are the connected components of
the complement of the walls.
\end{definition}

\begin{theorem}[Coupling polytopes are Delzant]\label{thm:delzant}
If $(f,g)$ is nondegenerate, then $\Cpl(f,g)$ is a Delzant polytope with respect to $\Lambda$: it is
simple, its facet normals are rational, and at each vertex the primitive edge vectors form a basis of
$\Lambda$. Its edges at a vertex are the fundamental cycles of a spanning tree.
\end{theorem}
\begin{proof}
Regard $\alpha$ as a weighting of the edges of the complete bipartite graph on $Y\sqcup Y'$.
\emph{Vertices are forests.} If the support of a feasible $\alpha$ contains a cycle, adding $\pm\epsilon$
alternately around it stays feasible, so $\alpha$ is not a vertex. If the support is a forest, the
marginal equations determine $\alpha$ on it uniquely (peel off leaves), so $\alpha$ is a vertex.
\emph{Nondegeneracy makes them spanning trees.} Every node has positive marginal and hence an incident
edge of the support. If the support forest had two or more components, take one with row set $I$ and
column set $J$. Then $\sum_If=\sum_Jg$; $I$ and $J$ are nonempty, and if $I=Y$ then $\sum_Jg=1$ forces
$J=Y'$ and the component is everything. So $I$ and $J$ are proper, contradicting nondegeneracy.
\emph{Simplicity and smoothness.} At a vertex with spanning tree $T$, the zero cells are the
$nm-(n+m-1)=(n-1)(m-1)$ non-tree cells $c$. Each closes a unique cycle in $T\cup\{c\}$, giving a
circulation $C_c$ with entries in $\{0,\pm1\}$ and $C_c(c)=1$. A circulation $w$ equals
$\sum_{c\notin T}w(c)C_c$, because the difference is a circulation supported on the tree, hence zero.
So the feasible directions at the vertex form the cone $\{\sum_ct_cC_c:t\ge0\}$ on exactly
$\dim\Cpl(f,g)$ independent generators (simplicity), and the $C_c$ are integral and form a $\mathbb{Z}$-basis
of $\Lambda$, since integer coefficients $w(c)$ suffice (smoothness). The facet normals are the coordinate
functionals $\alpha\mapsto\alpha(c)$, which are rational.
\end{proof}

By Delzant's theorem \cite{Delzant1988} there is a compact connected toric symplectic manifold
$(X_{f,g},\omega)$ of real dimension $2(n-1)(m-1)$, unique up to equivariant symplectomorphism, whose
moment map $\mu\colon X_{f,g}\to\Cpl(f,g)$ has image the coupling polytope. The fibre over an interior
coupling is a Lagrangian torus $T^{(n-1)(m-1)}$, and the fibre over a vertex is a point.

\begin{proposition}[Fisher metric and Guillemin's potential]\label{prop:guillemin}
Suppose every inequality $\alpha(c)\ge0$ defines a facet of $\Cpl(f,g)$. The canonical K\"ahler metric
of $X_{f,g}$ constructed by Guillemin \cite{Guillemin1994} has, in action--angle coordinates over the
interior, the form $\mathrm{Hess}\,G$ in the polytope directions, with symplectic potential
\[
G(\alpha)=\tfrac12\sum_{c}\alpha(c)\log\alpha(c).
\]
Hence the Fisher metric on $\Cpl(f,g)$ is twice the restriction of the canonical K\"ahler metric to the
real section, and the product coupling $f\otimes g$, the coupling of maximal entropy, is the unique
critical point of $G$ on the polytope.
\end{proposition}
\begin{proof}
Guillemin's potential is $\tfrac12\sum_i\ell_i\log\ell_i$ over the facets, where $\ell_i\ge0$ are the
facet inequalities written with primitive inward normals. Here $\ell_c(\alpha)=\alpha(c)$, and the
coordinate functional is primitive on $\Lambda$ because some integer circulation has entry $1$ at $c$.
Then $\mathrm{Hess}\,G(v,v)=\tfrac12\sum_cv(c)^2/\alpha(c)$. The critical-point condition is that
$\log\alpha$ is orthogonal to all circulations, that is $\log\alpha(y,y')=r(y)+s(y')$, which together with
the marginals gives $\alpha=f\otimes g$. Uniqueness follows from strict convexity.
\end{proof}

If some cell is not a facet, for instance when $f(y)+g(y')>1$ forces $\alpha(y,y')>0$, the potential
omits that term and the identity holds only up to those terms.

\subsection{Walls, chambers and monotone cases}

\begin{proposition}[Chambers]\label{prop:chambers}
The set of spanning trees that support vertices, and hence the combinatorial type of $\Cpl(f,g)$ and
the toric manifold $X_{f,g}$ up to equivariant diffeomorphism, is constant on each chamber. The type
can change across a wall. For $f=(\tfrac12,\tfrac12)$ and $g=(t,\tfrac23-t,\tfrac13)$ the polytope is a
hexagon for $t<\tfrac12$ and a quadrilateral for $t>\tfrac12$, both Delzant: two facets become redundant
across the wall $g_1=f_1$.
\end{proposition}
\begin{proof}
The tree solution on a spanning tree $T$ assigns to each tree edge the value $\sum_If-\sum_Jg$ (up to
sign) for the bipartition $(I,J)$ obtained by deleting that edge. So feasibility of $T$ depends only on
the signs of such differences, which are constant on chambers. The example is checked in
\texttt{check\_conditionals.py}.
\end{proof}

\begin{proposition}[Monotone cases]\label{prop:monotone}
Assume $(f,g)$ is nondegenerate and every cell defines a facet. Then $X_{f,g}$ is monotone if and only
if $f$ and $g$ are uniform, and uniform marginals are nondegenerate exactly when $\gcd(n,m)=1$. In
particular, in the square case $Y=Y'$ of coupling 2-cells no nondegenerate pair is monotone. On the
wall, the uniform $2\times2$ pair still gives a Delzant segment, $X=\mathbb{CP}^1$, with the product
coupling as its midpoint; the uniform $n\times n$ pair with $n\ge3$ gives a polytope that is not simple.
\end{proposition}
\begin{proof}
For a Delzant polytope written with primitive normals, the toric manifold is monotone exactly when
some interior point has equal lattice distance to all facets: the classes of the toric divisors span
$H^2$ subject to the linear relations given by the normals, $c_1$ is the sum of the divisor classes, and
$[\omega]$ is the sum weighted by the facet distances. Here the distances are the entries
$\alpha(c)$, so they are equal exactly when $\alpha$ is constant, which forces $f$ and $g$ to be
uniform; then $\alpha=f\otimes g$. Uniform marginals are nondegenerate iff $a/n\ne b/m$ for
$0<a<n$, $0<b<m$, iff $\gcd(n,m)=1$, which fails when $n=m$. The two wall cases are checked in
\texttt{check\_conditionals.py} (T3).
\end{proof}

\begin{example}[The degree-6 del Pezzo surface]\label{ex:dp6}
For $f$ uniform on two points and $g$ uniform on three, $\Cpl(f,g)$ is a hexagon. In the basis of
$\Lambda$ above, its facet normals are $\pm e_1$, $\pm e_2$ and $\pm(e_1+e_2)$. This is the fan of
$\mathbb{CP}^2$ blown up at three points, the del Pezzo surface of degree $6$. All six entries of the
product coupling equal $\tfrac16$, so it is the equidistant point. The case $3\times4$ (uniform) gives a
Delzant polytope of dimension $6$ with $96$ vertices and all $12$ cells facets, again monotone.
\end{example}

%=====================================================================
\section{The imported layer: Floer theory of the special fibre}\label{sec:imported}
%=====================================================================

We now quote results from symplectic topology. They are not proved or checked here.

\paragraph{Nondisplaceability.} For a monotone toric manifold, the torus fibre over the point at equal
distance from all facets (the \emph{special fibre}) is a stem in the sense of Entov and Polterovich,
hence nondisplaceable \cite{EntovPolterovich2006}. For toric Fano manifolds, Cho and Oh show that a
torus fibre with a flat $U(1)$ connection has nonvanishing Floer cohomology when it corresponds to a
critical point of the disc potential \cite{ChoOh2006}; see \cite{Cho2004} for the Clifford torus and
\cite{FOOO2010} for general compact toric manifolds.

\paragraph{Consequence for couplings.} In the monotone cases of \cref{prop:monotone}, the special fibre
lies over the product coupling $f\otimes g$. For \cref{ex:dp6}, the disc potential of the monotone
special fibre is, up to the factor $T^{1/6}$, the Laurent polynomial
\[
W(x,y)=x+y+\frac1x+\frac1y+xy+\frac1{xy},
\]
read off from the facet normals, and it has a critical point at $(x,y)=(1,1)$ (checked exactly), so the
special fibre with trivial holonomy has nonzero Floer cohomology by \cite{ChoOh2006}.

The maximal coupling, by contrast, has zero entries whenever $f\ne g$, so it lies on the boundary of
$\Cpl(f,g)$, where the torus fibre collapses (\texttt{check\_conditionals.py}, T6). The first draft's
conjecture that the critical points are the maximal couplings should therefore be replaced. The
geometrically distinguished coupling is the product coupling, which is also the maximum-entropy
coupling and the critical point of Guillemin's potential (\cref{prop:guillemin}).

%=====================================================================
\section{Corrections to the first draft}\label{sec:corrections}
%=====================================================================

\begin{enumerate}[leftmargin=*]
\item \emph{Candidate disc potential.} $W_\beta(f)=\sum_g\int_{\Cpl(f,g)}e^{-\beta c(\alpha)}\,d\nu$ sums over a
  continuum of 1-cells $g$ and is not defined as written. Its Laplace limit
  $\min_{g,\alpha}c(\alpha)$ is always $0$, attained at $g=f$, $\alpha=\Delta_f$, so it carries no
  information. The toric disc potential of $X_{f,g}$ is the correct object (\cref{sec:imported}).
\item \emph{Symplectic structure.} A symplectic form is antisymmetric and the Fisher metric is
  symmetric, so the restriction of the one cannot be the other. The correct statement is
  \cref{prop:guillemin}: the Fisher metric is twice the canonical K\"ahler metric on the real section
  of $X_{f,g}$.
\item \emph{Critical points.} ``Critical points of $\pi\colon E\to\Kl(\D)$'' is undefined, since $\pi$ is a
  functor, not a smooth map. The replacement is the special fibre over the product coupling in the
  monotone cases, and the walls of \cref{def:walls} as the locus where the family degenerates.
\item \emph{Monodromy and Dehn twists.} The family $(f,g)\mapsto X_{f,g}$ is constant on chambers
  and changes by toric birational modifications across walls (\cref{prop:chambers}); no vanishing
  cycle or Dehn twist has been identified. \Cref{conj:family} restates the question.
\item \emph{The $m_3$ conjecture.} Gluing is strictly associative (\cref{thm:witness}(i)), so each
  hom-category, regarded as an $A_\infty$-category, has $m_k=0$ for $k\ge3$. The whiskering defect
  compares $u\cdot(\beta\cdot\alpha)$ with $(u\cdot\beta)\cdot(u\cdot\alpha)$. Its inputs are two 2-cells and a
  1-cell, and it measures the failure of whiskering to be a functor. In $A_\infty$ language that is the
  second component $\mathcal{F}_2$ of an $A_\infty$-functor, not a ternary composition
  (\cref{conj:ainfty}).
\item \emph{Homotopy category.} The graded shadow has hom-sets $[d(f,g),\infty)$ under addition. These
  are not the degree-zero cohomology of cochain complexes, so the shadow is not a homotopy category of
  an $A_\infty$-category without further structure.
\item \emph{Fisher versus Fisher--Rao.} These are the same metric (\cref{prop:interior}).
\item \emph{References.} The companion tower is now version 9; its posted earlier version has the
  title ``\ldots and Weak Distributive Laws'', a claim withdrawn in the revisions.
\end{enumerate}

%=====================================================================
\section{The conjectural layer, restated}\label{sec:conjectural}
%=====================================================================

\begin{conjecture}[Family over the chamber complex]\label{conj:family}
The toric manifolds $X_{f,g}$ form a family over each chamber, and crossing a wall is a toric
birational transformation (a blow-up, blow-down or flip) determined by the pair $(I,J)$ defining the
wall. On a wall, $X_{f,g}$ degenerates to a singular toric variety whose moment polytope is the
non-simple $\Cpl(f,g)$.
\end{conjecture}

\begin{question}[Location of Floer-nontrivial fibres]\label{q:location}
Outside the monotone cases, the fibres with nonvanishing Floer cohomology (possibly after bulk
deformation) are located by the potential of \cite{FOOO2010}. Is the product coupling always among
them, and do they have a statistical meaning?
\end{question}

\begin{conjecture}[Whiskering as an $A_\infty$-functor]\label{conj:ainfty}
After a linearization of coupling 2-cells (for instance chains on the coupling polytopes), whiskering by
a 1-cell $u$ extends to an $A_\infty$-functor whose first component is the whiskering of
\cref{thm:witness}(ii) and whose second component is the whiskering defect. The $A_\infty$-functor
equation in arity two then reads
$\mathcal{F}_1(m_2(\beta,\alpha))-m_2(\mathcal{F}_1\beta,\mathcal{F}_1\alpha)=
m_1\mathcal{F}_2(\beta,\alpha)\pm\mathcal{F}_2(m_1\beta,\alpha)\pm\mathcal{F}_2(\beta,m_1\alpha)$.
Finding a linearization in which this holds is the first concrete target.
\end{conjecture}

\begin{conjecture}[Fukaya-type functor]\label{conj:fukaya}
There is a functor from a category built from the Fukaya categories of the $X_{f,g}$ (with special
fibres as distinguished objects) to a linearization of the coupling 2-cells, compatible with gluing on
one side and composition of Lagrangian correspondences on the other.
\end{conjecture}

%=====================================================================
\section{Next computations}\label{sec:next}
%=====================================================================

\begin{enumerate}[label=(\arabic*),leftmargin=*]
\item \emph{The $2\times3$ family.} Describe the toric surfaces in every chamber of
  $\{(f,g)\}$ for $\lvert Y\rvert=2$, $\lvert Y'\rvert=3$, and the transformations across each wall
  (\cref{conj:family}).
\item \emph{Square $2\times2$.} Here $\Cpl$ is a segment and $X=\mathbb{CP}^1$. Compute the coupling
  2-cells, the whiskerings and their defects exactly for all four deterministic 1-cells and a
  rational grid of stochastic ones, as the base case of \cref{conj:ainfty}.
\item \emph{Level two.} Extend \cref{thm:delzant} to couplings of second-order states, using the
  level-two gluing of \cite[Prop.~7.16]{Awasthi-Attribution}.
\end{enumerate}

\section{Computational checks}\label{sec:code}

All checks pass.
\begin{itemize}[leftmargin=*]
\item \texttt{check\_glue.py} (30 checks): \cref{prop:1cells,thm:witness}.
\item \texttt{check\_stochastic\_basechange.py} (13 checks): \cref{thm:fibre}; the hexagon of
  \cref{ex:dp6}; the Delzant property for random nondegenerate $3\times3$ marginals; non-simplicity of
  the uniform $3\times3$ case; the Hessian identity and critical point of \cref{prop:guillemin}.
\item \texttt{check\_conditionals.py} (7 checks): \cref{thm:delzant} by a geometric test (vertex degrees
  from the face lattice and unimodularity of primitive edge vectors) on random $2\times3$, $3\times3$
  and $2\times4$ cases; \cref{prop:chambers}; the degenerate $2\times2$ and $3\times3$ cases;
  \cref{prop:monotone,ex:dp6}, including the $3\times4$ case; the critical point of $W$; the boundary
  position of maximal couplings; strict associativity of gluing.
\end{itemize}

\begin{thebibliography}{99}

\bibitem{Awasthi-Attribution}
V.~K. Awasthi.
\newblock Precise attribution of model interventions: attribution sets, mechanism reach, path defects
and gluing.
\newblock Technical report, Johns Hopkins University, October 2026. Draft 4.

\bibitem{Awasthi-Tower}
V.~K. Awasthi.
\newblock The three-level posterior tower: hierarchical {B}ayesian inference and mixed distributive laws.
\newblock Technical report, Johns Hopkins University, October 2026. Revised version 9. An earlier
version is posted as DOI 10.13140/RG.2.2.29743.50088.

\bibitem{Cho2004}
C.-H. Cho.
\newblock Holomorphic discs, spin structures, and {F}loer cohomology of the {C}lifford torus.
\newblock \emph{International Mathematics Research Notices}, 2004(35):1803--1843, 2004.

\bibitem{ChoOh2006}
C.-H. Cho and Y.-G. Oh.
\newblock Floer cohomology and disc instantons of {L}agrangian torus fibers in {F}ano toric manifolds.
\newblock \emph{Asian Journal of Mathematics}, 10(4):773--814, 2006.

\bibitem{Delzant1988}
T.~Delzant.
\newblock Hamiltoniens p\'eriodiques et images convexes de l'application moment.
\newblock \emph{Bulletin de la Soci\'et\'e Math\'ematique de France}, 116(3):315--339, 1988.

\bibitem{EntovPolterovich2006}
M.~Entov and L.~Polterovich.
\newblock Quasi-states and symplectic intersections.
\newblock \emph{Commentarii Mathematici Helvetici}, 81(1):75--99, 2006.

\bibitem{FOOO2010}
K.~Fukaya, Y.-G. Oh, H.~Ohta, and K.~Ono.
\newblock Lagrangian {F}loer theory on compact toric manifolds {I}.
\newblock \emph{Duke Mathematical Journal}, 151(1):23--174, 2010.

\bibitem{Guillemin1994}
V.~Guillemin.
\newblock Kaehler structures on toric varieties.
\newblock \emph{Journal of Differential Geometry}, 40(2):285--309, 1994.

\end{thebibliography}
\end{document}
